\section*{Problem 1: Arithmetic in Prime Fields (Level 1)}

\subsubsection*{Mathematical part}
\textbf{Recall:} Before we can manipulate polynomials or construct extension fields, we need a notion of arithmetic where every operation remains inside a finite set. The simplest finite fields are the \emph{prime fields} $GF(p)$, where $p$ is a prime number.

Elements of $GF(p)$ are simply the integers $0, 1, \ldots, p-1$, and every arithmetic operation is performed modulo $p$. For example, in $GF(5)$:
\[
4+3 = 7 \equiv 2 \pmod 5
\]
Division requires finding the multiplicative inverse. Instead of searching for the inverse explicitly, we use Fermat's Little Theorem. If $p$ is prime and $a\neq0$, then $a^{p-1}\equiv1\pmod p$. Multiplying both sides by $a^{-1}$ gives:
\[
\frac{1}{a} = a^{p-2}
\]

\subsubsection*{Implementation part}
\textbf{Task:} Complete the \lstinline|GaloisFieldElement| class to support standard arithmetic via Python "dunder" (magic) methods. 

\textbf{Details:} From the programmer's perspective, every field element stores only two pieces of information: its numerical value, and the field to which it belongs. Your implementation must:
\begin{itemize}
    \item Perform the integer arithmetic.
    \item Reduce the result modulo the field's modulus.
    \item Return a \textbf{new} \lstinline|GaloisFieldElement| instance.
    \item Use binary exponentiation (exponentiation by squaring) for \lstinline|__pow__| to guarantee $O(\log n)$ complexity. Ensure $0^0 = 1$. You can assume the \lstinline|exponent| is always a non-negative integer.
\end{itemize}

\begin{center}
\renewcommand{\arraystretch}{1.5}
\begin{tabularx}{\textwidth}{@{} >{\bfseries}l X c @{}}
\toprule
Method & Arguments & Returns \\ \midrule
\_\_add\_\_ & \lstinline|self, other: GaloisFieldElement| & \lstinline|GaloisFieldElement| \\
\_\_sub\_\_ & \lstinline|self, other: GaloisFieldElement| & \lstinline|GaloisFieldElement| \\
\_\_mul\_\_ & \lstinline|self, other: GaloisFieldElement| & \lstinline|GaloisFieldElement| \\
\_\_truediv\_\_ & \lstinline|self, other: GaloisFieldElement| & \lstinline|GaloisFieldElement| \\
\_\_pow\_\_ & \lstinline|self, exponent: int| & \lstinline|GaloisFieldElement| \\
\bottomrule
\end{tabularx}
\end{center}

\subsubsection*{Experimentation part}
\textbf{UI Guide:} Open \lstinline|lab_dashboard.py| and navigate to the early debugging tabs if you wish to see how field elements wrap around algebraically.

\vspace{1cm}
